\begin{abstract}
A quadrotor deployed outside the lab meets conditions it was not tuned for: extra payload, a worn motor, wind and command latency. Meta-reinforcement learning promises adaptation from a few flights, while classical stacks built on geometric control, adaptive control, occupancy mapping and A* planning are dependable but fixed. We ask where meta-learning actually helps when it is combined with such a stack. In one simulator parameterised after a Holybro X500 V2, we insert learning at three points: a residual on the geometric controller, the controller's gains, and the local planner of a depth-camera, mapping and A* navigation stack. We compare seven algorithms (MAML, FOMAML, ANIL, Meta-SGD, Reptile, PEARL and RL$^2$) with nominal, L1-adaptive, domain-randomised and fine-tuning baselines on held-out and out-of-distribution dynamics. A learned residual cuts held-out tracking error from \valResLonePost\,cm with L1 adaptation to \valResDrPost\,cm, but gradient-based adaptation adds little beyond domain randomisation. Meta-learning helps most when it tunes classical gains (Meta-SGD: \valGainMetasgdPre\ to \valGainMetasgdPost\,cm). In navigation, gradient-based adaptation degrades learned planners while RL$^2$ improves with experience, and the classical planner remains the most reliable on the full stack in unseen rooms. Outside the training distribution every method fails, and learned residuals crash more often than the plain controller. We release the suite and recommend a hardware platform and staged plan for real-world tests.
\end{abstract}
